
\documentclass{article}
\usepackage{colortbl}
\usepackage{makecell}
\usepackage{multirow}
\usepackage{supertabular}

\begin{document}

\newcounter{utterance}

\centering \large Interaction Transcript for game `cladder', experiment `full\_v1.5\_default', episode 6482 with selfrefine{-}Qwen3.5{-}27B.
\vspace{24pt}

{ \footnotesize  \setcounter{utterance}{1}
\setlength{\tabcolsep}{0pt}
\begin{supertabular}{c@{$\;$}|p{.15\linewidth}@{}p{.15\linewidth}p{.15\linewidth}p{.15\linewidth}p{.15\linewidth}p{.15\linewidth}}
    \# & \multicolumn{2}{c}{Player} && \multicolumn{2}{c}{Game Master} \\
    \hline

    \theutterance \stepcounter{utterance}  
    & & & \multicolumn{4}{p{0.6\linewidth}}{
        \cellcolor[rgb]{0.9,0.9,0.9}{
            \makecell[{{p{\linewidth}}}]{
                \texttt{\tiny{[P1$\langle$GM]}}
                \texttt{You are an expert in causal inference. The following question is not a typical commonsense query, but rather a meticulously designed question created by a professor specializing in causal inference, intended to assess the students' mastery of the course content.} \\
\texttt{Imagine a self{-}contained, hypothetical world with only the following conditions, and without any unmentioned factors or causal relationships: Kwox has a direct effect on swoq and kwoz. Swoq has a direct effect on kwoz. For those who are not kwox, the probability of kwoz is 36\%. For those who are kwox, the probability of kwoz is 56\%. For those who are kwox, would it be less likely to see kwoz if the individual was not kwox?} \\
\texttt{Start your answer with 'Yes' or 'No', followed by additional reasoning or evidence to support your explanation.} \\
            }
        }
    }
     \\ \\

    \theutterance \stepcounter{utterance}  
    & & \multicolumn{4}{p{0.6\linewidth}}{
        \cellcolor[rgb]{0.95,0.95,0.95}{
            \makecell[{{p{\linewidth}}}]{
                \texttt{\tiny{[P1$\leftrightarrow$P]}}
                \texttt{{[}DRAFT{]}} \\
\texttt{Yes} \\
\\ 
\texttt{The question asks whether, for an individual who is currently **kwox**, the probability of observing **kwoz** would be lower if that individual were instead **not kwox**. This is a counterfactual query asking to compare the probability of the outcome under the current state (\$P(Kwoz | Kwox)\$) versus the probability under the alternative state (\$P(Kwoz | \textbackslash{}neg Kwox)\$).} \\
\\ 
\texttt{Based on the data provided in the hypothetical world:} \\
\texttt{1.  The probability of **kwoz** given **kwox** is **56\%** (\$P(Kwoz | Kwox) = 0.56\$).} \\
\texttt{2.  The probability of **kwoz** given **not kwox** is **36\%** (\$P(Kwoz | \textbackslash{}neg Kwox) = 0.36\$).} \\
\\ 
\texttt{The causal structure (Kwox \$\textbackslash{}to\$ Swoq \$\textbackslash{}to\$ Kwoz and Kwox \$\textbackslash{}to\$ Kwoz) explains *why* these probabilities differ, but the numerical values provided are sufficient to answer the comparison. Since 36\% is strictly less than 56\%, an individual who is **kwox** would indeed see a lower probability of **kwoz** (36\%) if they were hypothetically **not kwox**. Therefore, it is less likely to see kwoz in the counterfactual scenario where the individual is not kwox.} \\
            }
        }
    }
    & \\ \\

    \theutterance \stepcounter{utterance}  
    & & \multicolumn{4}{p{0.6\linewidth}}{
        \cellcolor[rgb]{0.95,0.95,0.95}{
            \makecell[{{p{\linewidth}}}]{
                \texttt{\tiny{[P1$\leftrightarrow$P]}}
                \texttt{{[}SELF{-}FEEDBACK{]}} \\
\texttt{FORMAT: ok} \\
\texttt{RULES: ok} \\
\texttt{STRATEGY: ok} \\
\texttt{REASON: The answer correctly starts with "Yes" as required, accurately interprets the counterfactual probability comparison using the provided data, and logically concludes that the probability decreases from 56\% to 36\%.} \\
            }
        }
    }
    & \\ \\

    \theutterance \stepcounter{utterance}  
    & \multicolumn{4}{p{0.6\linewidth}}{
        \cellcolor[rgb]{0.561,0.737,0.561}{
            \makecell[{{p{\linewidth}}}]{
                \texttt{\tiny{[P1$\rangle$GM]}}
                \texttt{Yes} \\
\\ 
\texttt{The question asks whether, for an individual who is currently **kwox**, the probability of observing **kwoz** would be lower if that individual were instead **not kwox**. This is a counterfactual query asking to compare the probability of the outcome under the current state (\$P(Kwoz | Kwox)\$) versus the probability under the alternative state (\$P(Kwoz | \textbackslash{}neg Kwox)\$).} \\
\\ 
\texttt{Based on the data provided in the hypothetical world:} \\
\texttt{1.  The probability of **kwoz** given **kwox** is **56\%** (\$P(Kwoz | Kwox) = 0.56\$).} \\
\texttt{2.  The probability of **kwoz** given **not kwox** is **36\%** (\$P(Kwoz | \textbackslash{}neg Kwox) = 0.36\$).} \\
\\ 
\texttt{The causal structure (Kwox \$\textbackslash{}to\$ Swoq \$\textbackslash{}to\$ Kwoz and Kwox \$\textbackslash{}to\$ Kwoz) explains *why* these probabilities differ, but the numerical values provided are sufficient to answer the comparison. Since 36\% is strictly less than 56\%, an individual who is **kwox** would indeed see a lower probability of **kwoz** (36\%) if they were hypothetically **not kwox**. Therefore, it is less likely to see kwoz in the counterfactual scenario where the individual is not kwox.} \\
            }
        }
    }
    & & \\ \\

    \theutterance \stepcounter{utterance}  
    & & & \multicolumn{2}{p{0.3\linewidth}}{
        \cellcolor[rgb]{0.9,0.9,0.9}{
            \makecell[{{p{\linewidth}}}]{
                \texttt{\tiny{[GM$|$GM]}}
                \texttt{Yes} \\
            }
        }
    }
    & & \\ \\

    \theutterance \stepcounter{utterance}  
    & & & \multicolumn{2}{p{0.3\linewidth}}{
        \cellcolor[rgb]{0.9,0.9,0.9}{
            \makecell[{{p{\linewidth}}}]{
                \texttt{\tiny{[GM$|$GM]}}
                \texttt{yes} \\
            }
        }
    }
    & & \\ \\

    \theutterance \stepcounter{utterance}  
    & & & \multicolumn{2}{p{0.3\linewidth}}{
        \cellcolor[rgb]{0.9,0.9,0.9}{
            \makecell[{{p{\linewidth}}}]{
                \texttt{\tiny{[GM$|$GM]}}
                \texttt{game\_result = WIN} \\
            }
        }
    }
    & & \\ \\

\end{supertabular}
}

\end{document}
